\section{The idea and the intuition behind it}
\label{sec:idea}

\subsection{Relations as nodes: the entity--connector graph}
\label{sec:idea-graph}

The graph has two classes of node (Figure~\ref{fig:capm}).

\termdef{Entity nodes} are the underlying quantities: the risk-free rate, the return of asset~\(i\), total current assets, the effective tax rate. An entity carries a canonical name and aliases, type metadata (kind of quantity, units, dimensionality), an optional constant value, a plain-language definition, and a vector embedding computed from that text for approximate matching.

\termdef{Connector nodes} are typed \(n\)-ary relations. A connector's \emph{kind} is an open registry; the first kind, the \emph{equation}, stores a symbolic expression in canonical \(\ldots = 0\) form, a binding map from the free symbols of the expression to the entities they stand for, per-symbol solvability metadata (closed form, closed form with a branch, numerical only, not solvable), and the assumptions under which the equation holds. Two further kinds were added during the year: the \emph{rule} (a stated constraint or inequality, never used to produce a value) and the \emph{aggregation} (a subtotal with one total and a signed, weighted part edge per member; Section~\ref{sec:impl-secgraph}).

\termdef{Edges} are undirected and minimal: two endpoints, a local symbol, a role. The one topology rule is that an edge never joins two entities; it always ends at a connector. A connector may, however, reference \emph{another connector} as a sub-relation, as the Black--Scholes price formula uses the definition of \(d_1\). The graph is therefore not strictly bipartite, which is what lets a complex formula be composed from named pieces.

Everything algebraic lives inside the connector's expression: products, exponents (including entity-valued exponents), repeated occurrences of one entity, and elementary functions such as \(\log\), \(\exp\) and the normal distribution function. An early design placed operators on the edges (``this edge carries \(E[\cdot]\)''); it could not express \(X^{Y}\) or \(X\cdot X\) and was removed within the first month. A second decision, taken in September after it blocked every CAPM question, is that \textbf{there are no statistical operators anywhere}: the expected value, variance or covariance of a quantity is a quantity of its own and therefore an entity of its own (``expected return of asset \(i\)''), bound by an ordinary symbol. The sympy-based parser now refuses an expression that applies \(E[\cdot]\) or \(\mathrm{Var}[\cdot]\), because an inert function can be neither inverted nor evaluated.

\begin{figure}[t]
\centering
\fitwidth{\begin{tikzpicture}[
  >=Stealth,
  ent/.style={circle,draw=cEntity,fill=cEntity!12,thick,minimum size=9mm,font=\scriptsize,align=center,inner sep=1pt},
  con/.style={rectangle,draw=cConnector,fill=cConnector!15,thick,rounded corners=1pt,minimum width=36mm,minimum height=11mm,font=\scriptsize,align=center,inner sep=3pt},
  lab/.style={font=\scriptsize\ttfamily,fill=white,inner sep=1pt},
  every edge/.style={draw,thick,gray!70}
]
% ---------- CAPM ----------
\node[con] (capm) at (0,0) {\textbf{equation} \\ $E_{R_i} - R_f - \beta_i\,(E_{R_m} - R_f) = 0$ \\ solvable for every symbol};
\node[ent,above left=11mm and 14mm of capm] (eri) {expected\\return\\of asset $i$};
\node[ent,above right=11mm and 14mm of capm] (rf) {risk-free\\rate};
\node[ent,below left=11mm and 14mm of capm] (beta) {beta of\\asset $i$};
\node[ent,below right=11mm and 14mm of capm] (erm) {expected\\market\\return};
\path (eri) edge node[lab,pos=.45]{E\_R\_i} (capm);
\path (rf) edge node[lab,pos=.45]{R\_f} (capm);
\path (beta) edge node[lab,pos=.45]{beta\_i} (capm);
\path (erm) edge node[lab,pos=.45]{E\_R\_m} (capm);
\node[font=\scriptsize,gray,align=center,below=18mm of capm] {$R_f$ appears twice in the body but is \emph{one} binding, one edge.\\ The realised return of asset $i$ is a different entity, which this equation does not touch.};

% ---------- connector referencing connector ----------
\begin{scope}[xshift=88mm]
\node[con] (d1) at (0,-14mm) {\textbf{equation} \texttt{d1\_def} \\ $d_1 - \dfrac{\ln(S/K)+(r+\sigma^2/2)T}{\sigma\sqrt T} = 0$};
\node[con] (bs) at (0,14mm) {\textbf{equation} \texttt{bs\_price} \\ $C - S\,\Phi(d_1) + K e^{-rT}\Phi(d_1-\sigma\sqrt T) = 0$};
\node[ent,left=9mm of d1] (S) {spot\\price $S$};
\node[ent,right=9mm of d1] (sig) {volatility\\$\sigma$};
\node[ent,left=9mm of bs] (C) {call\\price $C$};
\node[ent,right=9mm of bs] (K) {strike\\$K$};
\path (S) edge (d1); \path (sig) edge (d1); \path (C) edge (bs); \path (K) edge (bs);
\path (bs) edge[dashed,cConnector,very thick] node[lab,right=1pt,pos=.5]{d1 $\to$ \texttt{NodeRef(d1\_def)}} (d1);
\node[font=\scriptsize,gray,align=center,below=4mm of d1] {A connector may bind a symbol to another connector (dashed):\\ the graph is bipartite in spirit, not strictly.};
\end{scope}
\end{tikzpicture}}
\caption{The entity--connector graph. Left: CAPM as one equation connector (orange) bound to four entity nodes (blue) by undirected edges carrying the local symbol. The expected returns are entities in their own right; no edge or expression applies $E[\cdot]$. Right: a connector referencing another connector, the $d_1$ definition used inside the Black--Scholes call price. $r$ and $T$ are omitted from the drawing for space.}
\label{fig:capm}
\end{figure}


The intuition is that an equation is a thing in its own right, with a body and an identity, not a label on a pair of nodes. Once it is a node, it can own a symbolic body, declare its assumptions, be tested against data, be cited in a derivation, and be used in whichever direction a question requires.

\subsection{Solving as search}

A question names a target quantity at a coordinate (company, period, segment, \ldots). The solver treats the target as an OR-node over the connectors that can produce it and each connector as an AND-node over its other inputs, which are themselves targets. A value is found when a leaf is given in the question, a constant, or a figure read from a document. Three principles shape the search: \emph{trust before cost} (a derivation that uses only identities beats one touching an empirically fitted law, whatever the cost), \emph{cheapest data demand first} (among derivations of equal trust, the one needing the fewest looked-up quantities wins), and \emph{no silent wrong answers} (an alignment the solver cannot make, a unit it cannot carry, a cell holding several conflicting readings is reported as unsolved or put to a model, never guessed).

\subsection{The change of direction: from graph as store to graph as prior}
\label{sec:idea-pivot}

The original plan used the graph as the \emph{store}: read every table of a filing into facts, link each fact's label to a concept, stamp it with company, period, unit and scale, and at question time look the cell up and solve. This was built, measured on a real combined 10-K (seven registrants in one filing), and found to fail completely: zero of eleven questions, twice. The post-mortem measured each ingestion decision in isolation (Table~\ref{tab:commitment}). None was terrible; all were taken \emph{once}, blind to any question, and an answer needed every one of them to be right for every fact it touched. The product of those rates is well under one half before the solver runs.

\begin{table}[h]
\centering\small
\caption{Accuracy of each decision taken at ingestion, measured on the Entergy 2011 combined 10-K (September 2026). Each is made once, with no question in view; an answer needs all of them.}
\label{tab:commitment}
\begin{tabular}{lr}
\toprule
decision & measured \\
\midrule
row label linked to the right concept & 79\% of facts \\
period correct & 93\% of timed facts \\
company correct & \(\approx\)84\% of facts carrying one \\
cell holds exactly one distinct value & 69\% of linked cells \\
\bottomrule
\end{tabular}
\end{table}

The conclusion was that the failure is not the graph but \emph{commitment at ingest}. The opposite policy, late binding, stores nothing but the text and its vectors and makes every decision at question time, when the question supplies the context each decision was missing: ``Entergy Texas'' and ``2010 to 2011'' say which table to read and which column to trust. A plain text-and-tables retrieval system was built as the control (Section~\ref{sec:impl-nograph}) and immediately answered questions the structured path had refused.

That left the question of what the graph is \emph{for}. The answer that organises the rest of the year is that the graph is a \textbf{prior}: background knowledge that shapes how the text is searched, read and checked, rather than the place where facts live. Concretely, the graph supplies (i)~the formula and its inputs when a quantity is not printed, (ii)~the vocabulary a question's words map onto, (iii)~the arithmetic, recomputed in code from cited inputs, and, through two further knowledge bases, (iv)~\emph{where} to look. Documents, meanwhile, are read by a model at question time under one hard rule: a figure enters an answer only if the code finds its quote printed in the cited passage, in the cited row, with the sign and the unit the page gives.

\subsection{Three knowledge bases}
\label{sec:idea-three}

The analyst's second skill, knowing where a quantity is printed, turned out to need knowledge the equation graph does not hold. Two more bases were therefore designed and built in November (Figure~\ref{fig:three}):

\begin{description}
\item[The concept graph] (the entity--connector graph above): \emph{what} quantities exist and how they relate. Built from textbooks and from the SEC's reporting taxonomy, with every relation verified against filings before admission.
\item[The source graph]: \emph{kinds of documents} and their official parts. A 10-K has Items 1--15 in four parts; Item~8 holds the statements; a 10-Q covers a quarter and a year to date and never a fourth quarter; an 8-K with Item~2.02 attaches an earnings release; a 10-K/A amends a 10-K. Each entry is hand-authored from the official SEC text with a verified quote, and carries the periods its parts hold in a closed vocabulary of about thirteen relative words (own period end, prior year end, year to date, quarters ahead, \ldots).
\item[The bridge]: \emph{where each quantity is printed} --- in which part and statement or note of which kind of document, for which periods, how often --- measured over 32,288 filings from the SEC's structured data sets, supplemented by what the system observes while answering.
\end{description}

\begin{figure}[t]
\centering
\fitwidth{\begin{tikzpicture}[
  >=Stealth, font=\small,
  base/.style={draw,thick,rounded corners=3pt,align=left,inner sep=6pt,text width=46mm,font=\scriptsize},
  kb/.style={base,fill=cLight},
  title/.style={font=\small\bfseries},
  arr/.style={->,thick,gray!80}
]
\node[kb,draw=cEntity] (cg) at (0,0) {%
 \textcolor{cEntity}{\textbf{Concept graph}} --- \emph{what} and \emph{how related}\\[2pt]
 9{,}121 entities (4{,}474 textbook, 4{,}647 SEC)\\
 5{,}457 relations: 2{,}753 equations, 1{,}384 subtotals, 1{,}320 rules\\
 4{,}983 links entity $\leftrightarrow$ SEC concept, one entity per concept\\[2pt]
 \emph{verified by}: arithmetic over filings; review of definitions\\
 \emph{example}: current ratio $=$ current assets / current liabilities (identity)};
\node[kb,draw=cDoc,right=14mm of cg] (sg) {%
 \textcolor{cDoc}{\textbf{Source graph}} --- \emph{kinds of documents}\\[2pt]
 15 kinds, 150 official parts, 29 period rules, 20 relations, 4 filer classes\\
 closed vocabulary of $\sim$13 period words\\[2pt]
 \emph{verified by}: a quote from the official SEC text per entry; counts over 68{,}177 filings\\
 \emph{example}: 10-K $\triangleright$ Part II $\triangleright$ Item 8 holds the statements for the fiscal year and the prior year end; a 10-Q never holds Q4};
\node[kb,draw=cConnector,below=12mm of $(cg.south)!0.5!(sg.south)$] (br) {%
 \textcolor{cConnector}{\textbf{Bridge}} --- \emph{where each quantity is printed, how often}\\[2pt]
 49{,}677 rows (concept, kind, part, statement or note, periods, count, of)\\
 measured over 32{,}288 filings of 7 kinds from the SEC Financial Statement (and Notes) data sets; observed rows added from the system's own finds; authored rows as priors\\[2pt]
 \emph{example}: effective tax rate $\to$ 10-K, Item 8, income-tax note, 56\% of filings; capex $\to$ cash flow statement, 61\%};
\draw[arr] (cg.south) -- ++(0,-5mm) -| ([xshift=-12mm]br.north);
\draw[arr] (sg.south) -- ++(0,-5mm) -| ([xshift=12mm]br.north);
\node[font=\scriptsize,gray,align=center,below=3mm of br] {The bridge is keyed on the SEC concept (stable) and resolved to the entity through the concept links when opened.};
\end{tikzpicture}}
\caption{The three knowledge bases. The concept graph says what quantities exist and how they relate; the source graph says what a document of a given kind contains and for which periods; the bridge joins them by saying in which part and statement of which kind each quantity is printed, with counts. Each is verified by a different kind of evidence.}
\label{fig:three}
\end{figure}


The intuition for the division is that the three bases change at different rates and are verified by different evidence. Equations are timeless and are tested by arithmetic over filings. Document kinds are fixed by regulation and are tested by quotation. Where a quantity is printed is a statistical fact about filers' behaviour and is \emph{counted}. Keeping them apart keeps each one honest and lets the question-answering loop consult the one it needs at each step.

\subsection{Learned where learning is safe, authored where it is not}

One further principle settled during the year. Pieces whose behaviour can be fitted from data the system is never evaluated on are fitted, not hand-tuned: the retriever's score over kinds of evidence is a conditional logit trained on EDGAR filings that label their own numbers through inline XBRL, with the benchmark companies excluded; the probabilities it reports, and the bridge's counts, are calibrated against held-out filings. Pieces that are judgements about meaning (does this filer's caption denote this concept? is this cited cell the target?) are made by a model at question time, every time, with no cache, because a judgement made for one question's context has selected the wrong fact in another. And pieces that assume the \emph{shape} of a document (identity tie-outs, range checks on a value) are not enforced in code at all: a check that a citation is printed stays hard; a check against the graph asks the model once and records a disagreement as a flag.
